\scanpage{103}

Lorsque $k$ est de caractéristique nulle, on a toujours un foncteur fibre naturel: le \uline{foncteur de Hodge,} qui à la cohomologie motivique d'une variété (projective lisse) $X$ associe le vectoriel bigradué $\handwritten{\coprod}H^q(X,\handwritten{\underline{\Omega}}^p_{X/k})$ Donc pour le groupe de Galois motivique correspondant $G$, on trouve un homomorphisme naturel
\[
\underline{G}_m^2\handwritten{\longrightarrow}G,
\]
i.e. deux homomorphismes commutant l'un à l'autre
\[
i_1,i_2:\underline{G}_m\handwritten{\longrightarrow}G.
\]
Le fait que la graduation totale dans la cohomologie de Hodge corresponde au poids des motifs s'exprime par la relation
\[
i(\handwritten{\lambda})=i_1(\handwritten{\lambda})\,i_2(\handwritten{\lambda});
\]
on fera attention que $i_1$ et $i_2$ ne sont pas centraux (car la bigraduation en cohomologie de Hodge ne correspond pas à une bigraduation d'un motif !)

2) Prenant toujours pour $\underline{C}$ la catégorie des motifs sur $k$, avec $k$ de car. nulle, on a un autre foncteur fibre canonique, le \uline{foncteur de De Rham} qui associe à la cohomologie motivique d'une variété $X$ le vectoriel $H^*(X,\handwritten{\Omega}^*_{X/k})$ (espace d'hypercohomologie). Ce vectoriel n'est plus bigradué\handwritten{\unsure{é,|en,}} mais seulement gradué et filtré,\handwritten{(}la filtration étant celle associée à la suite spectrale d'hypercohomologie, commençant avec la cohomologie de Hodge\handwritten{)}. On sait d'ailleurs que cette suite spectrale dégénère (théorie de Hodge), donc $\operatorname{Gr}(H_{\mathrm{DR}}(V))\handwritten{\simeq}H_{\mathrm{Hdg}}(V)$ (isomorphisme fonctoriel en le motif $V$\inserted{\handwritten{où on \unsure{gradue} le $2^{\mathrm{ème}}$ membre \uline{par $p$}}}). On peut se proposer d'analyser à quelle structure supplémentaire, sur le groupe de Galois motivique associé au foncteur fibre $H_{\mathrm{DR}}$, correspond la filtration canonique de ce foncteur.
% The inserted handwritten note is circled, with a pointer immediately after V.

De façon générale, étant donné une catégorie tensorielle $C$ sur $k$ munie d'un foncteur fibre $F$, on peut se proposer de déterminer les filtra%
% archive sheet number: 102
